% Generated draft: requires mathematical and source review.
\documentclass[11pt]{article}
\usepackage[margin=25mm]{geometry}
\usepackage{fontspec}
\setmainfont{Latin Modern Roman}
\usepackage{amsmath,amssymb,mathtools,unicode-math}
\setmathfont{Latin Modern Math}
\usepackage{enumitem,needspace,xurl,hyperref}
\setlength{\parindent}{0pt}
\setlength{\parskip}{5pt}
\setlength{\emergencystretch}{3em}
\newcommand{\N}{\mathbb{N}}
\newcommand{\Z}{\mathbb{Z}}
\newcommand{\Q}{\mathbb{Q}}
\newcommand{\R}{\mathbb{R}}
\newcommand{\C}{\mathbb{C}}
\newcommand{\F}{\mathbb{F}}
\newcommand{\A}{\mathbb{A}}
\newcommand{\PP}{\mathbb{P}}
\newcommand{\E}{\mathbb{E}}
\newcommand{\eps}{\varepsilon}
\newcommand{\dd}{\,\mathrm{d}}
\begin{document}
\section*{UnsolvedMath ID: 30000874}
\textbf{Source problem code:} catalogue entry
\subsection{Definitions and mathematical statement}
This is an empirical mathematics-education research question, rather than a conjecture asserting a mathematical implication. Its target is teachers' mathematical knowledge both as content they need to know and as knowledge they need to know how to use. The source asks:
\begin{quote}
What distinct approaches help teachers develop the mathematical knowledge they need to know and know how to use, and what evidence shows how these approaches function and with what effects?
\end{quote}
A source-faithful account must therefore identify and distinguish professional-development approaches, address the development of both mathematical knowledge and its usable application, and state the evidence for how each approach functions and what effects are observed. The supplied assessment additionally emphasizes the conditions under which an approach works. These are linked descriptive, explanatory, and evaluative tasks, not an assertion that one intervention succeeds for every teacher or context.

The record does not operationally define the required knowledge, its use, distinctness of approaches, outcome measures, comparison conditions, teacher population, timescale, or standard of evidence. It supplies no dataset or quantified effect statement. Consequently, no probability model, causal estimand, universal effectiveness theorem, or prescribed study design can be introduced as if it were part of the source. A precise empirical formulation requires those missing research choices, while retaining both knowledge and use as targets.

The supplied assessment dated 2026-08-21 says that distinct approaches and outcome studies exist, but reports no general evidence-based classification of which approaches work, under what conditions, and with what effects. This is supplied literature triage, not a verified claim that all existing studies are inadequate or that no particular approach has supporting evidence.
\subsection{Short English statement}
Identify approaches that develop teachers' mathematical knowledge and ability to use it, and explain the evidence for their mechanisms, conditions, and effects. The source poses an empirical research program and supplies no fixed measurements, population, or quantitative theorem.
\subsection{Sources}
\begin{itemize}
\item Local UnsolvedMath registry, permanent ID \(30000874\), exact source code \texttt{catalogue entry}, proposed year 2007; no individual proposer is supplied.
\item \emph{Professional Development of Mathematics Teachers---Research and Practice from an International Perspective (2008)}, Oberwolfach Reports, DOI 10.4171/owr/2007/52, pp. 3073--3130. \url{https://doi.org/10.4171/owr/2007/52}.
\item Supplied reference labelled ``international perspective review'': \url{https://doi.org/10.1016/j.jmathb.2018.05.002}.
\item Supplied reference labelled ``systematic teacher-knowledge review'': \url{https://doi.org/10.1007/s10857-025-09737-8}.
\item The record labels the statement \texttt{exact} and reports checking its content against the official report. The open label and assessment checked 2026-08-21 are dated supplied triage requiring independent verification of the review scope and claims.
\end{itemize}
\end{document}
